% PAGES: 53-54

\noindent The numerical solution of linear systems.
\medskip

\noindent Forward substitution. Backward substitution.
\medskip

\noindent Finding discount factors using forward substitution.
\medskip

\noindent LU decomposition without pivoting. Existence and uniqueness. Pseudocode and
operation count.
\medskip

\noindent Linear solvers using the LU decomposition without pivoting.
\medskip

\noindent LU linear solvers for tridiagonal matrices.
\medskip

\noindent LU decomposition with row pivoting. Pseudocode and operation count.
\medskip

\noindent Linear solvers using the LU decomposition with row pivoting.
\medskip

\noindent Solving linear systems corresponding to the same matrix.
\medskip

\noindent Finding discount factors using the LU decomposition.
\medskip

\noindent Cubic spline interpolation. Cubic spline interpolation for zero rate curves.
\medskip

\section{The numerical solution of linear systems}

A fundamental problem in numerical analysis is solving linear systems, i.e., finding
a vector $x$ such that
\begin{equation}
Ax \;=\; b,
\end{equation}
where $A$ is a square matrix and $b$ is a column vector of size equal to the number of
columns of $A$.

The \textit{numerical solution} of the problem (2.1) requires finding efficient
numerical algorithms to compute the vector $x$ such that $Ax = b$, given the matrix $A$
and the column vector $b$. These numerical algorithms will also address implicitly the
question whether the matrix $A$ is nonsingular, which is required for the linear system
(2.1) to have a unique solution.

The numerical methods for solving linear systems are of two fundamentally different
types:

\begin{itemize}
\item \textit{Direct Methods:} The solution of the linear system $Ax=b$ is found by
computing the LU decomposition with pivoting of the matrix $A$, or the Cholesky
decomposition of $A$, if the matrix $A$ is symmetric positive definite, and then
solving two linear systems corresponding to lower triangular matrices and to upper
triangular matrices.
\item \textit{Iterative Methods:} The solution of the linear system $Ax=b$ is obtained
recursively. The iteration is stopped once an approximate solution with prescribed
precision is found. Examples of such methods are Jacobi iteration, Gauss--Siedel
iteration, and SOR (Successive Overrelaxation).
\end{itemize}

All the methods mentioned above are much more efficient from a numerical standpoint
than a computation of the inverse matrix $A^{-1}$ and finding $x$ by using the
matrix--vector multiplication $x=A^{-1}b$.

In practice, direct methods are faster for solving linear systems arising from
discretizations of two dimensional problems, while iterative methods are faster for
solving linear systems corresponding to three dimensional problems (or higher
dimensional problems). We only discuss direct methods herein: see sections 2.4 and 2.6
for the LU decomposition, and section 6.1 for the Cholesky decomposition.

Note that linear systems corresponding to nonsingular diagonal matrices have
straightforward solution. If $D$ is a nonsingular diagonal matrix of size $n$, i.e.,
with $D(j,j)\neq 0$ for all $j=1:n$, see Lemma 1.12, then the solution $x$ to $Dx=b$,
where $b$ is an $n\times 1$ vector, is given by
\[
x(j) \;=\; \frac{b(j)}{D(j,j)}, \quad \forall\, j=1:n.
\]

Linear systems corresponding to lower triangular matrices or to upper triangular
matrices can also be solved directly, by using the Forward Substitution, if the matrix
is lower triangular, or by using the Backward Substitution, if the matrix is upper
triangular; see section 2.2 and section 2.3, respectively.
